\begin{center}\begin{tikzpicture}[font=\small,>=stealth,line join=round]
\draw[gray!40,rounded corners=3pt] (0,0) rectangle (14.2,8.3);\node at (7.1,7.700000000000001) {例5：整段取2、整段取3，走两段回到原值};
\begin{scope}[shift={(1.7315789473684213,2.45)},x=0.7894736842105263cm,y=1.0cm]\draw[->,gray] (-0.8,0)--(6.8,0) node[right]{$x$};\draw[blue,thick] (0,2)--(2,2);\draw[blue,fill=white] (2,2) circle(2.7pt);\fill[blue] (0,2) circle(2.3pt);\node[below=7pt] at (0,-0.75){$0$};\draw[orange!85!black,thick] (2,3)--(4,3);\draw[orange!85!black,fill=white] (4,3) circle(2.7pt);\fill[orange!85!black] (2,3) circle(2.3pt);\node[below=7pt] at (2,-0.75){$2$};\draw[blue,thick] (4,2)--(6,2);\draw[blue,fill=white] (6,2) circle(2.7pt);\fill[blue] (4,2) circle(2.3pt);\node[below=7pt] at (4,-0.75){$4$};\fill[orange!85!black] (6,3) circle(2.3pt);\node[below=7pt] at (6,-0.75){$6$};\end{scope}\node[anchor=west,align=left,text width=5cm] at (8.4,4){每段长度2\\[9pt]$\Phi(u)=\frac{u+1}{u-1}$\\[9pt]$2\mapsto3\mapsto2$\\[9pt]$f(0)=2,\ f(2)=3$\\[9pt]$2026$ 与2同余于周期4};
\node[text width=12.6cm,align=center] at (7.1,.65) {每段左闭右开；完整函数在实数轴上按长度4重复。};
\end{tikzpicture}\end{center}
